\documentclass[12pt]{article}
\usepackage[utf8]{inputenc}
\usepackage[T1]{fontenc}
\usepackage[margin=1in]{geometry}
\usepackage{setspace}
\usepackage{xcolor}
\usepackage[colorlinks=true, urlcolor=blue!60!black, linkcolor=blue!60!black]{hyperref}
\setlength{\parindent}{2em}
\setlength{\parskip}{1.25em}
\renewcommand{\baselinestretch}{1.0}
\title{\textit{Ziglar v. Abbasi}, 582 U.S. 120 (2017)}
\author{Case Brief}
\date{}
\begin{document}
\maketitle

\section{Case Metadata}

\begin{itemize}
  \item \textbf{Case Name:} \textit{Ziglar v. Abbasi} (consolidated with \textit{Ashcroft v. Abbasi}, No. 15--1359, and \textit{Hasty v. Abbasi}, No. 15--1363)
  \item \textbf{Citation:} 582 U.S. 120 (2017) (\href{https://www.law.cornell.edu/supremecourt/text/15-1358}{Cornell LII Primary Opinion}; \href{https://supreme.justia.com/cases/federal/us/582/15-1358/}{Justia}; \href{https://www.supremecourt.gov/opinions/16pdf/15-1358_6khn.pdf}{SupremeCourt.gov Slip Opinion})
  \item \textbf{Court:} Supreme Court of the United States (on writs of certiorari to the U.S. Court of Appeals for the Second Circuit, \textit{Turkmen v. Hasty}, 789 F.3d 218 (2d Cir. 2015))
  \item \textbf{Dates:} Argued January 18, 2017; Decided June 19, 2017
  \item \textbf{Vote:} 4--2 (reversing in part, vacating and remanding in part)
  \item \textbf{Majority Author:} Justice Anthony M. Kennedy (opinion of the Court as to Parts I, II, III, IV--A, and V, joined by Chief Justice Roberts and Justices Thomas and Alito; plurality opinion as to Part IV--B, joined by Chief Justice Roberts and Justice Alito)
  \item \textbf{Concurring Opinion:} Justice Thomas (concurring in part and concurring in the judgment)
  \item \textbf{Dissenting Opinion:} Justice Breyer (joined by Justice Ginsburg)
  \item \textbf{Recusals:} Justices Sotomayor, Kagan, and Gorsuch took no part in the consideration or decision of the cases
  \item \textbf{Chapter Placement:} Chapter 1: Judicial Review and the Federal Courts (Implied Constitutional Causes of Action under \textit{Bivens}; Separation of Powers; Qualified Immunity)
\end{itemize}

\section{Facts and Procedural History}

In the immediate aftermath of the September 11, 2001 terrorist attacks, the Federal Bureau of Investigation (FBI) investigated more than 96,000 public tips and arrested over 700 noncitizens on immigration violations. Detainees designated ``of interest'' to the terrorism investigation---or whose status remained uncertain---were held without bail under a ``hold-until-cleared'' policy (\href{https://www.law.cornell.edu/supremecourt/text/15-1358}{582 U.S. at 126--127}). Respondents, six men of Arab or South Asian descent (five of whom were Muslim) present in the United States without legal authorization, were among roughly 84 detainees confined in the Administrative Maximum Special Housing Unit of the Metropolitan Detention Center (MDC) in Brooklyn, New York, for periods ranging from three to eight months before release and removal (\textit{id.} at 127--128).

Respondents alleged that official Bureau of Prisons policy kept them in tiny cells over 23 hours a day with lights on 24 hours, denied basic hygiene items and outside communications, required four-guard shackled escorts, and mandated frequent strip searches. They further alleged unofficial ``physical and verbal abuse'' by guards---including slamming detainees into walls, breaking bones, and religious taunts---that MDC wardens knowingly tolerated (\textit{id.} at 127--128).

Respondents brought a putative class action in the Eastern District of New York against two groups of federal officials: (1) ``Executive Officials''---former Attorney General John Ashcroft, former FBI Director Robert Mueller, and former Immigration and Naturalization Service Commissioner James Ziglar; and (2) ``Wardens''---MDC Warden Dennis Hasty and Associate Warden James Sherman. Invoking \href{https://supreme.justia.com/cases/federal/us/403/388/}{\textit{Bivens v. Six Unknown Named Agents of Federal Bureau of Narcotics}}, 403 U.S. 388 (1971), they asserted four constitutional damages claims: Fifth Amendment substantive due process and equal protection challenges to the detention policy against all petitioners, Fourth and Fifth Amendment strip-search claims against the Wardens, and a Fifth Amendment deliberate-indifference prisoner-abuse claim against the Wardens. They also alleged an equal-protection conspiracy under 42 U.S.C. \S~1985(3). After \textit{Ashcroft v. Iqbal}, 556 U.S. 662 (2009), the District Court dismissed the claims against the Executive Officials but allowed claims against the Wardens to proceed. A divided Second Circuit panel (\textit{Turkmen v. Hasty}, 789 F.3d 218 (2015), Judge Raggi dissenting) reinstated the claims against the Executive Officials and affirmed most claims against the Wardens (dismissing only the abuse claim against Sherman), and denied rehearing en banc (808 F.3d 197 (2015)). The Supreme Court granted certiorari (580 U.S. 915 (2016)).

\section{Issue(s) Presented}

\begin{enumerate}
  \item Under what analytical framework may federal courts imply a \textit{Bivens} damages remedy directly under the Constitution, and does a \textit{Bivens} remedy extend to (a) high-level post-9/11 detention-policy claims or (b) a Fifth Amendment deliberate-indifference prisoner-abuse claim against a federal warden?
  \item Are federal officials within the same executive department entitled to qualified immunity from a civil-conspiracy damages claim under 42 U.S.C. \S~1985(3)?
\end{enumerate}

\section{Holding}

\textbf{Reversed in part; vacated and remanded in part (4--2).}

\begin{itemize}
  \item \textbf{Formalized Two-Step \textit{Bivens} Framework:} Expanding \textit{Bivens} is a ``disfavored'' judicial activity. Courts must first ask whether a claim arises in a ``new \textit{Bivens} context''---meaning it differs in any meaningful way from the Court's three recognized \textit{Bivens} precedents (\textit{Bivens}, \textit{Davis v. Passman}, and \textit{Carlson v. Green}). If so, courts may not imply a remedy if ``special factors counselling hesitation'' indicate Congress is better suited to weigh the costs and benefits of a damages action (\href{https://www.law.cornell.edu/supremecourt/text/15-1358}{582 U.S. at 135--140}).
  \item \textbf{Detention-Policy Claims Barred:} No \textit{Bivens} remedy lies for respondents' detention-policy claims against either the Executive Officials or the Wardens; the claims arose in a new context and were foreclosed by national-security, executive-deliberation, congressional-silence, and alternative-remedy special factors (\textit{id.} at 140--146).
  \item \textbf{Prisoner-Abuse Claim Against Warden Hasty Vacated and Remanded:} Although respondents pleaded a plausible Fifth Amendment violation, the claim still presented a ``new context'' distinct from \textit{Carlson}, requiring vacatur and remand for a special-factors analysis (\textit{id.} at 146--149).
  \item \textbf{Qualified Immunity Under \S~1985(3):} All petitioners were entitled to qualified immunity because it was not clearly established whether officers within the same federal department can ``conspire'' with one another under 42 U.S.C. \S~1985(3) (\textit{id.} at 150--156).
\end{itemize}

\section{Reasoning}

\subsection{The Shift Away from Implied Remedies and the Two-Step Test (Parts II--III)}

Justice Kennedy framed \textit{Bivens} (Fourth Amendment search and seizure), \href{https://supreme.justia.com/cases/federal/us/442/228/}{\textit{Davis v. Passman}}, 442 U.S. 228 (1979) (Fifth Amendment congressional sex discrimination), and \textit{Carlson v. Green}, 446 U.S. 14 (1980) (Eighth Amendment failure to provide prison medical care) as products of an ``ancien regime'' in which courts routinely implied private rights of action to effectuate statutory and constitutional purposes (\href{https://www.law.cornell.edu/supremecourt/text/15-1358}{582 U.S. at 131--132}). Although preserving \textit{Bivens} as settled law in its core domestic law-enforcement sphere, Kennedy emphasized that modern separation-of-powers doctrine places the decision ``who should decide'' whether to create damages liability---Congress or the courts---squarely on Congress (\textit{id.} at 134--136).

The Court rejected the Second Circuit's permissive ``same right and same mechanism of injury'' test and announced a strict two-step inquiry:
\begin{enumerate}
  \item \textbf{Step One (New Context):} A context is ``new'' whenever a case differs ``in a meaningful way'' from \textit{Bivens}, \textit{Davis}, or \textit{Carlson}. Non-exhaustive indicators include the rank of the officers involved, the constitutional right at issue, the generality or specificity of official action, the extent of judicial guidance, the officer's statutory mandate, the risk of disruptive judicial intrusion into coordinate branches, or potential special factors not previously considered (\textit{id.} at 139--140).
  \item \textbf{Step Two (Special Factors):} If the context is new, a court must refrain from implying a remedy whenever ``sound reasons'' suggest Congress might doubt the efficacy or necessity of a damages action, or when an ``alternative remedial structure'' exists (\textit{id.} at 136--137).
\end{enumerate}

Applying this framework to the detention-policy claims, the Court identified four decisive special factors: (1) \textit{Bivens} is not a vehicle to alter high-level executive policy, and discovery into senior DOJ deliberations would chill candid advice and distract officials; (2) challenging the Government's post-9/11 response directly implicates national-security policy committed to Congress and the President; (3) Congress's nearly 16-year silence---despite active oversight and a 300-page DOJ Inspector General report on MDC conditions---suggested intentional refusal to authorize damages; and (4) detainees had potential alternative remedies, including injunctive relief and habeas corpus (\textit{id.} at 140--145).

\subsection{Warden Hasty and \S~1985(3) Qualified Immunity (Parts IV--V)}

On the guard-abuse claim against Warden Hasty, the Court agreed that respondents plausibly alleged deliberate indifference, but stressed that ``even a modest extension is still an extension'' (\textit{id.} at 147). Because the claim rested on the Fifth rather than Eighth Amendment, involved less-settled supervisory standards, implicated potential equitable remedies, and followed Congress's omission of a standalone damages remedy in the Prison Litigation Reform Act of 1995 (PLRA), the ``new-context inquiry is easily satisfied'' (\textit{id.} at 148--149). A three-Justice plurality (Part IV--B) vacated and remanded for a special-factors analysis.

Finally, on the 42 U.S.C. \S~1985(3) conspiracy claim, the Court held petitioners were entitled to qualified immunity. Because circuits were sharply split over whether the antitrust-derived intracorporate-conspiracy doctrine (\textit{Copperweld Corp. v. Independence Tube Corp.}) bars conspiracy liability among agents of the same governmental department---and because open intra-departmental policy consultation should not be chilled---reasonable officials lacked clearly established notice that their joint policy discussions violated \S~1985(3) (\textit{id.} at 152--156).

\section{Concurring and Dissenting Opinions}

\begin{itemize}
  \item \textbf{Justice Thomas (Concurring in Part and in the Judgment):} Joined Parts I, II, III, IV--A, and V, but refused to join Part IV--B's remand. Quoting Justice Scalia's \textit{Malesko} concurrence, Thomas described \textit{Bivens} as ``a relic of the heady days in which this Court assumed common-law powers to create causes of action'' and would have limited \textit{Bivens} to its exact facts and reversed outright as to Hasty, concurring in the remand judgment solely to produce a controlling disposition (\href{https://www.law.cornell.edu/supremecourt/text/15-1358}{582 U.S. at 157}). Separately, in an influential critique, Thomas urged the Court in an appropriate case to reconsider modern ``clearly established'' qualified immunity because it departs from 1871 common-law tort immunities and reflects judicial ``freewheeling policy choice[s]'' (\textit{id.} at 157--160).
  \item \textbf{Justice Breyer (Dissenting, joined by Justice Ginsburg):} Argued the claims arose in familiar \textit{Bivens} and \textit{Carlson} territory---civilian detainees alleging physical abuse, strip searches, and invidious discrimination by federal officers---and warned the majority was shrinking, not merely declining to extend, \textit{Bivens}. Breyer contended that Congress had ratified \textit{Bivens} through the Westfall Act (28 U.S.C. \S~2679(b)(2)(A)), that communications blackouts rendered habeas or injunctive relief illusory (leaving ``damages or nothing''), and that pleading standards (\textit{Iqbal}), qualified immunity, and tailored discovery already protected executive officials during emergencies without abolishing post-crisis damages accountability (\textit{id.} at 160--183).
\end{itemize}

\section{Key Quotations}

\begin{quote}
``Given the notable change in the Court's approach to recognizing implied causes of action, however, the Court has made clear that expanding the \textit{Bivens} remedy is now a `disfavored' judicial activity.''\\
--- Justice Kennedy, \href{https://www.law.cornell.edu/supremecourt/text/15-1358}{582 U.S. at 135}
\end{quote}

\begin{quote}
``The proper test for determining whether a case presents a new \textit{Bivens} context is as follows. If the case is different in a meaningful way from previous \textit{Bivens} cases decided by this Court, then the context is new.''\\
--- Justice Kennedy, \href{https://www.law.cornell.edu/supremecourt/text/15-1358}{582 U.S. at 139}
\end{quote}

\begin{quote}
``The necessary inference, though, is that the inquiry must concentrate on whether the Judiciary is well suited, absent congressional action or instruction, to consider and weigh the costs and benefits of allowing a damages action to proceed. Thus, to be a `special factor counselling hesitation,' a factor must cause a court to hesitate before answering that question in the affirmative.''\\
--- Justice Kennedy, \href{https://www.law.cornell.edu/supremecourt/text/15-1358}{582 U.S. at 136}
\end{quote}

\begin{quote}
``Given these safeguards against undue interference by the Judiciary in times of war or national-security emergency, the Court's abolition, or limitation of, \textit{Bivens} actions goes too far. If you are cold, put on a sweater, perhaps an overcoat, perhaps also turn up the heat, but do not set fire to the house.''\\
--- Justice Breyer, dissenting, \href{https://www.law.cornell.edu/supremecourt/text/15-1358}{582 U.S. at 181}
\end{quote}

\section{Doctrinal Significance and Modern Legacy (Through 2026)}

\textit{Ziglar v. Abbasi} transformed \textit{Bivens} from a general constitutional-tort doctrine into a near-categorical barrier against implied damages suits against federal officers, setting off a decade-long doctrinal contraction continuing through 2026:

\begin{itemize}
  \item \textbf{Tightening the \textit{Abbasi} Vise (\textit{Hernandez}, \textit{Egbert}, and \textit{Goldey}):} In \href{https://supreme.justia.com/cases/federal/us/589/17-1678/}{\textit{Hernandez v. Mesa}}, 589 U.S. 93 (2020), the Court applied \textit{Abbasi} to foreclose a cross-border shooting claim against a Border Patrol agent. In \href{https://supreme.justia.com/cases/federal/us/596/21-147/}{\textit{Egbert v. Boule}}, 596 U.S. 482 (2022), the Court compressed \textit{Abbasi}'s two steps into a single overriding question---``whether there is any reason to think that Congress might be better equipped to create a damages remedy''---barring Fourth Amendment excessive-force and First Amendment retaliation claims against a Border Patrol officer. Most recently, in \href{https://supreme.justia.com/cases/federal/us/606/24-809/}{\textit{Goldey v. Fields}}, 606 U.S. 942 (2025) (per curiam) (\href{https://www.supremecourt.gov/opinions/24pdf/24-809_9o6b.pdf}{slip op.}), the Court summarily reversed the Fourth Circuit, citing \textit{Abbasi} (582 U.S. at 135, 137, 148--149) to hold that no \textit{Bivens} remedy exists for an Eighth Amendment excessive-force claim against federal prison guards given PLRA silence and internal BOP grievance procedures.
  \item \textbf{Broader Separation-of-Powers Reach Through 2026:} \textit{Abbasi}'s separation-of-powers skepticism toward judge-made causes of action now anchors adjacent federal-courts doctrine, exemplified by \href{https://www.supremecourt.gov/opinions/25pdf/24-856_kjfm.pdf}{\textit{Cisco Systems, Inc. v. Doe}}, 609 U.S. \_\_\_ (June 23, 2026) (foreclosing new judicially created Alien Tort Statute claims). Meanwhile, lower-court splits over what remains of \textit{Carlson} continue to reach the Court in 2026 (e.g., \textit{Nielsen v. Watanabe}, No. 25--417), fueling renewed legislative proposals to enact a federal statutory counterpart to 42 U.S.C. \S~1983.
  \item \textbf{Qualified Immunity Debate:} Justice Thomas's concurrence in Part II catalyzed a major cross-ideological originalist challenge to the ``clearly established'' qualified-immunity standard that remains prominent in scholarship and litigation through 2026.
\end{itemize}

\section{Pedagogical Critiques and Exam Hypotheticals}

\subsection{Critiques}

\begin{itemize}
  \item \textbf{The Rights-Without-Remedies Anomaly:} Critics echo Justice Breyer's concern that \textit{Abbasi} magnifies an asymmetric remedial gap: state officers face statutory damages under 42 U.S.C. \S~1983 for identical constitutional violations, while federal officers increasingly escape personal damages liability altogether.
  \item \textbf{Illusory ``Alternative Remedies'':} Treating prospective injunctions, unsettled habeas review of confinement conditions, or nonbinding agency grievance procedures as adequate alternatives ignores that post-release or communications-blackout detainees face a practical ``damages or nothing'' reality.
  \item \textbf{Unstated Overruling in Slow Motion:} Because ``the new-context inquiry is easily satisfied'' and any plausible policy concern counsels hesitation, commentators argue \textit{Abbasi} and \textit{Egbert} functionally confine \textit{Bivens} to its exact 1971 facts without candidly overruling it.
\end{itemize}

\subsection{``Opposite Side'' Contrast: \textit{Bivens} and \textit{Snowden v. Henning}}

To test the narrow outer boundary of what still survives \textit{Abbasi}, contrast \textit{Abbasi} with \href{https://supreme.justia.com/cases/federal/us/403/388/}{\textit{Bivens}}, 403 U.S. 388 (1971), and the Seventh Circuit's post-\textit{Egbert} application in \href{https://media.ca7.uscourts.gov/cgi-bin/OpinionsWeb/processWebInputExternal.pl?Submit=Display&Path=Y2023/D06-27/C:21-1463:J:Sykes:aut:T:fnOp:N:3066917:S:0}{\textit{Snowden v. Henning}}, 72 F.4th 237 (7th Cir. 2023). In \textit{Snowden}, a DEA agent allegedly used excessive force while executing a domestic arrest warrant at a hotel. Applying \textit{Abbasi}'s ``meaningful difference'' factors, the Seventh Circuit allowed the Fourth Amendment claim to proceed because a line-level narcotics officer making a routine domestic arrest implicated neither national security, high-level policy, nor a new category of defendants---demonstrating that only claims closely tracking \textit{Bivens}'s core domestic law-enforcement facts avoid Step One dismissal.

\subsection{Exam Hypotheticals}

\begin{enumerate}
  \item \textbf{Line Officer vs. Policy Official:} Suppose (a) a federal ATF agent breaks down a citizen's front door without a warrant during a domestic firearms investigation and punches the homeowner, while (b) the DHS Secretary issues an emergency directive ordering warrantless electronic sweeps of airports after a cyber-threat. Under \textit{Abbasi}, claim (a) closely parallels \textit{Bivens} (line-level officer, domestic Fourth Amendment search/seizure, clear judicial guidance; cf. \textit{Snowden}), though defendants will still argue even minor differences (agency or statutory mandate) create a new context under \textit{Egbert}. Claim (b) is squarely barred under \textit{Abbasi}: high-level executive policy, national security, and prospective injunctive alternatives all counsel hesitation.
  \item \textbf{Federal Pretrial Abuse After \textit{Goldey}:} Suppose a federal Marshal beats a pretrial detainee in a courthouse holding cell and denies him insulin for 48 hours. Students should separate the two theories under \textit{Abbasi} and \textit{Goldey v. Fields} (2025): the excessive-force claim fails as a new context foreclosed by special factors, while the medical-indifference claim tests whether the Fifth Amendment (pretrial detainee) vs. Eighth Amendment (convicted prisoner in \textit{Carlson}) distinction alone creates a ``new context'' under \textit{Abbasi} Part IV--A.
\end{enumerate}

\end{document}
